\subsection{Combustion Experiments and Diagnostics}

\subsubsection{What roles do canonical flames play in combustion research?}

{\color{gray}\footnotesize [Claude context: check Lecture 5 (Experimental Combustion 1 -- premixed vs.\ diffusion flames, Bunsen burner); Maas notes Ch.2 / Warnatz Ch.2 (Experimental Investigation of Flames)]}

\subsubsection{Laminar canonical flames are provided on standardized burners, for instance: the laminar premixed flame (Bunsen burner), the laminar ``rich'' premixed flame (McKenna burner), and the laminar diffusion flame (Wolfhard--Parker burner). Describe their laboratory flame geometries and the working principle of these burners. How can these burners be used to study combustion?}

{\color{gray}\footnotesize [Claude context: check Lecture 5 (Bunsen burner, premixed/diffusion flames); Lecture 9 (premixed/diffusion flame structure); Maas notes Ch.2 / Warnatz Ch.2]}

\subsubsection{What are the three main advantages of optical diagnostics in flames?}

{\color{gray}\footnotesize [Claude context: check Lecture 6 (Experimental Combustion 2 -- diagnostics) and Lecture 7 (Laser Diagnostics); Maas notes Ch.2 / Warnatz Ch.2]}

\subsubsection{What is temporal-, spatial-, and spectral resolution/control?}

{\color{gray}\footnotesize [Claude context: check Lecture 7 (laser-based diagnostic techniques); Maas notes Ch.2 / Warnatz Ch.2]}

\subsubsection{What is the Boltzmann thermal population distribution and how can it be used for thermometry in combustion?}

{\color{gray}\footnotesize [Claude context: check Lecture 7 (light-matter interaction, LIF); Maas notes Ch.2 / Warnatz Ch.2 (temperature measurements)]}

\subsubsection{How can species be detected and identified in flames? Why do we need this?}

{\color{gray}\footnotesize [Claude context: check Lecture 7 (laser diagnostics, scattering/LIF); Maas notes Ch.2 / Warnatz Ch.2 (concentration measurements)]}

\subsubsection{Give three examples of light--matter interaction and describe their working principles.}

{\color{gray}\footnotesize [Claude context: check Lecture 7 (light-matter interaction -- absorption, excitation, scattering, LIF); Maas notes Ch.2 / Warnatz Ch.2]}

\subsubsection{If you were given the assignment of measuring the laminar burning velocity of the simplest hydrocarbon fuel, methane, how would you do it?}

{\color{gray}\footnotesize [Claude context: check Lecture 5 (measuring laminar flame speed with Bunsen burner); Maas notes Ch.2 / Warnatz Ch.2 (Experimental Investigation of Flames); Maas notes Ch.7 / Warnatz Ch.8.3 (flame velocities)]}
